\section{Candidate Selection vis E-Graph Extraction}\label{sec:beam}

In this section we formalize \emph{beam extraction}, our approach
to selecting an optimal subset of abstractions
from a wider set of candidate abstractions in the presence of an e-graph.

\subsection{Library Selection as E-Graph Extraction}

Having generated potential candidates for learned abstractions
through e-graph anti-unification, the problem of LLMT now becomes
\emph{candidate selection}: which subset of learned abstractions
should we use to optimally rewrite our target term?
%
In the presence of e-graphs, this answer to this question is
non-trivial. As mentioned in \autoref{sec:egraphs}, not only can
$\denot{\mathcal{G}}$ be potentially infinite, making enumerating terms in
$\mathcal{G}$ impractical, but the choice of which subset of abstractions
is optimal depends on the concrete term itself.

To tackle this problem, we reduce the problem of library selection to
that of extracting a single, optimal concrete term from an e-graph.
%
Similarly to how we define a correspondence between patterns and
compression rules, which rewrite an expression into an application
of an abstraction,
we can also define a correspondence between patterns and
\emph{library rules}, which rewrite an expression into a combined
abstraction definition and application.
%
For a given \T{let}-binding $\T{let}\ P = \lambda \many{X} \to p$,
let its library rule $\libpat(P, \many{X}, p)$ be the rewrite rule
$p \Rightarrow \T{let}\ P = \lambda \many{X} \to p\ \T{in}\ P\ \many{X}$.
%
As an example, $\libpat(P, X, X + 1)$ corresponds to the rewrite rule
$X + 1\Rightarrow \T{let}\ P = \lambda X \to X + 1\ \T{in}\ P\ X$.
%
By generating a rewrite $\libpat(P_i, \fv(p_i), p_i)$ for every pattern
$p_i$ generated during e-graph anti-unification of an e-graph $\mathcal{G}$,
and applying these rewrites via equality saturation to $\mathcal{G}$, we
obtain an e-graph $\mathcal{G}'$ in which every possible application to a new
learned abstraction is included in the e-graph, always accompanied by an associated
abstraction definition at every application.
%
Because every application of a learned abstraction in any subterm of
any $t \in \denot{\mathcal{G}'}$ must be accompanied by a corresponding
\T{let}-binding, the subset of abstractions used in $t$ can be determined
solely by its \T{let}-bindings. Thus, finding the optimal set of libraries
to use for $\mathcal{G}'$ can be found by extracting the smallest term in
$\mathcal{G}'$, where ``smallest'' is given by a cost function $\cs$
which only counts the definitions of \T{let}-bindings once.
%
The set of abstractions which are \T{let}-bound in this term will be
the optimal subset of libraries to use for the programs represented by
$\mathcal{G}'$.

% \TODO{Now that we have all these abstraction definitions in our e-graph,
% every term that we can extract. We know that if an abstraction is used, it
% MUST be let-bound, the only way we can introduce an abstraction is through
% $\libpat$, which also introduces a let binding. Thus, every term in this
% new e-graph corresponds to a set of libraries based on what's let-bound in the
% term. From here, finding the best set of libraries is just e-graph term extraction:
% we find the smallest term, given that we only count let definitions once, then
% look at what's let bound. And that's our answer!}

\subsection{Optimal Extraction with Partial Shared Structure}

Having reduced library subset selection to the problem of extracting
a term from an e-graph in the presence of shared structure, the problem
now becomes defining a cost function $\cs$ which can simultaneously
account for common \T{let}-defined sub-expressions while also allowing
for efficient optimal term extraction.
%
Traditional e-graph term extraction won't work in this context, since
this cost function cannot be \TODO{not monotonic but another word
goes here}, and greedy term extraction
inherently relies on having a \TODO{same} cost function.
%
For instance, consider the following e-graph with eight e-classes,
corresponding to the term $f(g(a, 1, a), g(b, 1, a), g(c, 1, c))$
after the library rewrite $\libpat(l, x, g(x, 1, x))$ has been applied:
%
\begin{align*}
c_0&\colon \{ a \},\ c_1\colon \{ b \},\ c_2\colon \{ c \},\ c_3\colon \{ 1 \} \\
c_4&\colon \{ g(c_0, c_3, c_0), \T{let}\ l = \lambda x \to g(x, 1, x)\ \T{in}\ l(c_0) \} \\
c_5&\colon \{ g(c_1, c_3, c_1), \T{let}\ l = \lambda x \to g(x, 1, x)\ \T{in}\ l(c_1) \} \\
c_6&\colon \{ g(c_2, c_3, c_2), \T{let}\ l = \lambda x \to g(x, 1, x)\ \T{in}\ l(c_2) \} \\
c_7&\colon \{ f(c_4, c_5, c_6) \} \\
\end{align*}
%
A purely greedy extraction approach based on AST size would dictate that
$g(a, 1, a)$ should be the selected subterm at $c_4$, since it has the smallest
AST size of all the expressions in this e-class.
%
However, when used in the context of $c_7$, if the definition of $f$
is shared, then $\T{let}\ l = \lambda x \to g(x, 1, x)\ \T{in}\ l(a)$
should be selected, since the cost of the extra \T{let}-binding
can be amortized by its later use in $c_7$.
%
Indeed, in general, greedy extraction would \emph{never} extract a
use of an abstraction, since the cost of defining and applying
an abstraction always exceeds the cost of \TODO{the expression
itself} at the point of definition;
this cost only becomes \TODO{worth it} if the abstraction is
used at differet parts of the e-graph, which \TODO{a greedy
extraction algorithm wouldn't be able to surmise}.

Because our common sub-expression extraction problem is
inherently limited---only definitions in a \T{let}-binding are
considered shared---we can define a cost function which keeps track
of \T{let}-defined symbols and counts them once, while counting all
other subterms without regard for common sub-expressions.
%
By separating common sub-expression cost from \TODO{regular node?}
cost, we can \TODO{find the true cost of a term by finding
the sum of libs and expr cost.}

\begin{figure}
  \textbf{E-node cost set} $\cs_N(n)$
  $$
  \begin{array}{rl}
    \cs_N(s()) &= \{ (\emptyset, 1) \} \\
    \cs_N(s(a)) &= \{ (l, c + 1) \mid (l, c) \in \cs(a) \} \\
    \cs_N(s(\many{a_i})) &= \mathsf{cross}(\many{\cs(a_i)}) \\
    \cs_N(\T{let}\ s = \lambda \many{X} \to a\ \T{in}\ b) &= \{ (l, c + 1) \mid (l, c) \in \mathsf{clib}(s, \cs(a), \cs(b)) \} \\
    \end{array}
  $$

  \textbf{E-class cost set} $\cs(a)$
  $$
  \begin{array}{rl}
    \cs(\{ \many{n_j} \}) &= \bigcup \many{\cs_N(n_j)} \\
  \end{array}
  $$

  \textbf{Argument cost set combination} $\mathsf{cross}$
  $$
  \begin{array}{rl}
    \mathsf{cross}(z_1, z_2) &= \{ (l_1 \cup l_2, c_1 + c_2) \mid (l_1, c_1) \in z_1, (l_2, c_2) \in z_2 \} \\
  \end{array}
  $$

  \textbf{Abstraction cost set combination} $\mathsf{clib}$
  $$
  \begin{array}{rl}
    \mathsf{clib}(s, z_1, z_2) &= \{ (l_1 \cup l_2 \cup \{ (s, c_1) \}, c_2) \mid (l_1, c_1) \in z_1, (l_2, c_2) \in z_2 \} \\
  \end{array}
  $$
\caption{\TODO{caption for costset def}}\label{fig:costset}
\end{figure}

\autoref{fig:costset} describes the cost function calculation for
e-nodes and e-classes.
%
Every e-node and e-class has a corresponding \emph{cost set},
which represents a set of possible library selections for all terms
represented by the e-node or e-class:
a choice of a set of abstractions, their associated cost, and the
size of the rewritten term using these abstractions.
%
These cost sets are propagated bottom-up through the e-graph.

To justify our cost function definition, we revisit our previous
example, \TODO{in which greedy extraction would not have worked}.
%
For the e-class $c_0$, there is only one possible selection
of abstractions and one possible rewritten term cost for all the
terms in each e-class: the empty set of abstractions with term cost $1$.
%
This cost set, $\{ (\emptyset, 1) \}$, is the same as that of $c_1$,
$c_2$, and $c_3$.
%
To find the cost set for $c_4$, we must first find the cost set
calculations for each of its e-nodes.
%
Finding the cost set of the expression $g(c_0, c_3, c_0)$ requires
us to combine the cost sets of the argument e-classes $c_0$
and $c_3$;
in our formalization, this is performed by the $\mathsf{cross}$
operation.
%
The possible library subsets which can be used to rewrite $g(c_0, c_3, c_0)$
consists of any subset which can be used for $c_0$ combined
with any subset which can be used for $c_3$; the corresponding
term cost in each case is simply the term costs of $c_0$ and $c_3$ added
together (since this outer expression includes the subterms for $c_0$
and $c_3$) and incremented, to account for the $+$ AST node.
%
In this case, since the cost sets of $c_0$ and $c_3$ only contain one
element each, the cost set of $c_0 + c_3$ consists of one element:
$\{ (\emptyset \cup \emptyset, 1 + 1 + 1)\} = \{ (\emptyset, 3) \}$.
%
This corresponds to the intuition that the term $a + 1$ can be
expressed with the empty set of libraries, in which case it has term
size $3$.
%
For an e-node with three or more arguments, the $\mathsf{cross}$
operation is performed as a fold over the arguments of the e-node.

Finding the cost set for the other e-node in $c_4$,
$\T{let}\ f = \lambda x \to x + 1\ \T{in}\ f(c_0)$,
requires first finding the cost set for the application, $f(c_0)$.
%
The cost set of a unary e-node is equal to the cost set of
its argument, where all term costs have been incremented,
so this application would have a cost set of
$\{ (\emptyset, 2) \}$.
%
To find the cost set of this expression as a whole, we must
add the abstraction $f$ and its associated cost into every
possibility represented by the cost set of the \T{let}-binding's
body.
%
One complication with this is that the definition of the
abstraction itself is an e-class which can include applications
to other learned abstractions, and thus has its own cost set.
%
For the sake of simplicity, we assume the body of $f$ contains
one possible term, $\lambda x \to x + 1$, which would have
a cost set of $\{ (\emptyset, 3) \}$.
%
In this case, we see that the cost set of this expression should
be $\{ (\{ (f, 3) \}, 2) \}$, representing the fact that
this e-node can be expressed with term size 2 given an abstraction
$f$ of term size 3.
%
In general, this is handled by the $\mathsf{clib}$ operation,
which performs a similar function as $\mathsf{cross}$: for every
possible definition of an abstraction (and thus every possible
abstraction subset and term size it can have), it must be
combined with every possible abstraction and term size of the
\T{let}-binding's body.
%
In this case, ``combined'' means that all abstraction definitions,
including the \T{let}-bound abstraction itself and any other
abstractions it might use, are \TODO{unioned} together,
while the term sizes of the body are kept constant.
%
With the cost sets of these two e-nodes calculated, we can now
find the cost set for the e-class $c_4$ as a whole.
%
The cost set of an e-class is defined as the union of the cost sets
of its e-nodes;
the possible selections of library subsets in an e-class is equal to
all of the possible selections of library subsets of any of its
constituent e-nodes.
%
In this case, this means that
$\cs(c_4) = \{ (\emptyset, 3), (\{ (f, 3) \}, 2) \}\}$.
%
In other words, a term $t \in \denot{c_4}$ can have AST
size 3 if rewritten with no abstractions, or term size 2
if rewritten with the abstraction $f$ of size 3..
%
The cost sets of $c_5$ and $c_6$ are the same as $c_4$, and
the cost set for $c_7$ can be found in a similar fashion,
resulting in a final cost set:

\begin{align*}
\cs(c_7) = \{ &(\emptyset, 10), (\{ (f, 3) \}, 9), (\{ (f, 3) \}, 9), (\{ (f, 3) \}, 9), \\
&(\{ (f, 3) \}, 8), (\{ (f, 3) \}, 8), (\{ (f, 3) \}, 8), (\{ (f, 3) \}, 7) \}
\end{align*}

Library subset extraction is performed by examining the root
e-class and finding the library subset selection which results
in the lowest \emph{full cost}: the sum of the costs of the
shared expressions, the abstraction definitions, combined with
the term cost.
%
In this case, \TODO{that lowest is given by this.
Thus, the optimal library is this. Woo!}

\mypara{Partial Order Reduction}
%
The size of the cost set of an e-node is exponential with respect
to the number of arguments it has.
\TODO{This is not very performant. But we can prune a bunch of redundant
things. Consider the cost set from this last example, in which we define
a lib but never actually use it. we can remove this possibility completely!}

\mypara{From Optimal to Beam Extraction}
%
Even in the presence of partial order reduction, calculating this cost
set for every e-node and e-class in an e-graph representing a large
program can still be prohibitively expensive.
%
In practice, \tool provides multiple parameters which limit the number
of candidate libraries at each node (the \emph{beam size}),
as well as the size of each library.